\documentclass[11pt,a4paper]{article}

\usepackage[margin=1in]{geometry}
\usepackage{amsmath, amssymb, amsthm}
\usepackage{mathtools}
\usepackage{bm}
\usepackage{booktabs}
\usepackage{enumitem}
\usepackage{hyperref}
\usepackage{algorithm}
\usepackage{algpseudocode}

\newcommand{\bc}{\bm{c}}
\newcommand{\bx}{\bm{x}}
\newcommand{\by}{\bm{y}}
\newcommand{\bs}{\bm{s}}
\newcommand{\bb}{\bm{b}}
\newcommand{\E}{\mathbb{E}}
\newcommand{\R}{\mathbb{R}}
\newcommand{\argmax}{\operatorname*{arg\,max}}
\newcommand{\Var}{\operatorname{Var}}

\newtheorem{lemma}{Lemma}

\title{Stochastic Batched Dynamic Programming\\
       \large A scaling strategy for joint hardware-policy co-design beyond the exact-DP horizon ceiling}
\author{}
\date{}

\begin{document}
\maketitle

\begin{abstract}
Exact Bellman dynamic programming for partially observable Markov decision processes (POMDPs)
admits an efficient sufficient-statistic memoization for many problems of interest, but the
memoization grows fast enough in horizon $K$ that exact solutions are typically capped at
$K \approx 4$--$5$.  Joint hardware-policy co-design, where a continuous hardware parameter
$\bc$ is optimized jointly with the inner Bellman-optimal policy $\pi^\star(\bc)$, inherits this
ceiling.  We develop a scaling extension: \emph{stochastic batched dynamic programming},
which combines (i) a batched receding-horizon decomposition of the inner Bellman problem
into $n$ tractable batches of length $K_\mathrm{batch}$ and (ii) stochastic Monte-Carlo
sampling of batch-boundary trajectories at outer-gradient time, replacing the intractable
enumeration over reachable batch-boundary beliefs.  The resulting estimator is unbiased,
per-sample gradients are exact via the sharp-max envelope theorem, and the wall-clock cost
scales linearly in the sample size and the number of batches rather than in the reachable
belief count.  The construction generalizes naturally and preserves the exactness
guarantees of joint dynamic programming on the sampled sub-problems.
\end{abstract}

\tableofcontents

\section{Setup and notation}
\label{sec:setup}

A POMDP is a tuple $(\mathcal{X}, \mathcal{S}, \mathcal{Y}, p_0, p, T)$ with hidden state
$\bx \in \mathcal{X}$, action $\bs \in \mathcal{S}$, observation $\by \in \mathcal{Y}$,
prior $p_0$, observation likelihood $p(\by \mid \bx, \bs;\, \bc)$, and Bayesian belief
update operator $T(\bb, \bs, \by;\, \bc)$.  The hardware parameter $\bc \in \mathcal{C}$
parametrises the likelihood and the belief update.  Belief at step $k$ is
$\bb_k \in \Delta(\mathcal{X})$, the simplex of distributions on the state space.

A pointwise reward $\Phi(\bx, \bs_{1:K};\, \bc)$ depends on the realized trajectory and the
hardware.  For terminal rewards expressible as a functional of the terminal belief
$\bb_K$, write $\Phi_0(\bb_K;\, \bc) = \E_{\bx \mid \bb_K}[\Phi(\bx, \bs_{1:K};\, \bc)]$.
The Bellman recursion at fixed $\bc$ is
\begin{align}
V_K(\bb;\, \bc) &= \Phi_0(\bb;\, \bc), \label{eq:terminal}\\
V_k(\bb;\, \bc) &= \max_{\bs \in \mathcal{S}}\; \E_{\by \mid \bb, \bs, \bc}\!\bigl[V_{k+1}(T(\bb, \bs, \by;\, \bc);\, \bc)\bigr], \label{eq:bellman}
\end{align}
for $k = K-1, \ldots, 0$.  The optimal adaptive policy is
$\pi^\star_k(\bb;\, \bc) = \argmax_{\bs} \E_{\by}[V_{k+1}(T(\bb, \bs, \by;\, \bc);\, \bc)]$.
The joint hardware-policy optimum solves
\begin{equation}
\bc^\star = \argmax_{\bc \in \mathcal{C}}\; V^\star(\bc), \qquad V^\star(\bc) := V_0(p_0;\, \bc).
\label{eq:joint}
\end{equation}

The outer gradient $dV^\star/d\bc$ is computed by the sharp-max envelope theorem: at the
optimal inner policy, the indirect contribution of policy sensitivity to $\partial V^\star /
\partial \bc$ vanishes, and reverse-mode automatic differentiation through an
argmax-frozen Bellman recursion returns an exact, unbiased outer gradient.

\section{The exact-DP horizon ceiling}
\label{sec:ceiling}

Exact Bellman DP for the recursion of Equations~\eqref{eq:terminal}--\eqref{eq:bellman}
admits efficient memoization when a low-dimensional sufficient statistic for the belief is
available.  Two common settings:

\paragraph{Discrete state.}  The belief is a probability vector on $|\mathcal{X}|$ states.
Reachable beliefs from $p_0$ form a tree of size at most $(|\mathcal{S}| \cdot
|\mathcal{Y}|)^K$.  Memoization stores one value per reachable belief; the table size is
bounded but grows exponentially with $K$.

\paragraph{Continuous state with sufficient-statistic reduction.}  When the per-shot
likelihood depends on the history only through a finite tuple of integers (e.g., a
count-tuple of observations partitioned by action type), the belief at step $k$ is
encoded by that tuple, and the reachable-tuple set is the memo's natural index.  The
number of reachable tuples grows multiplicatively with $K$, polynomially with the action
and observation alphabet sizes.  In favourable cases the memo is $\sim 10^7$ at $K = 4$;
unfavourable cases (multimodal posteriors, dense action grids) push this to $10^8$--$10^9$
or more by $K = 5$ or $6$.  Memory becomes the binding constraint.

For practical hardware (modest workstation, $\sim 1$~TB RAM) the exact-DP ceiling is
typically $K_\mathrm{max} \in \{4, 5\}$.  Beyond that we descend to approximations.

\section{Batched dynamic programming}
\label{sec:batched}

Suppose the deployment horizon $K_\mathrm{total}$ is much larger than the exact-DP ceiling
$K_\mathrm{batch} \le K_\mathrm{max}$, with $K_\mathrm{total} = n \cdot K_\mathrm{batch}$
for some integer $n$.  Define a \emph{batched policy} $\pi^\mathrm{batch}$ as follows:

\begin{enumerate}[nosep, leftmargin=2em]
\item At step $k = 0$, with belief $\bb_0 = p_0$, solve exact Bellman for $K_\mathrm{batch}$
      epochs starting at $\bb_0$.  Use the resulting policy $\pi^{(1)}$ for steps $1, \ldots, K_\mathrm{batch}$.
\item At step $k = K_\mathrm{batch}$, with belief $\bb_{K_\mathrm{batch}}$ resulting from
      executing $\pi^{(1)}$ on the realized observations, solve exact Bellman for $K_\mathrm{batch}$
      epochs starting at $\bb_{K_\mathrm{batch}}$.  Use $\pi^{(2)}$ for the next $K_\mathrm{batch}$ steps.
\item Repeat for $n$ batches total.
\end{enumerate}

Each batch is locally Bellman-optimal over its $K_\mathrm{batch}$-step subproblem; the
overall procedure is a receding-horizon approximation to the full $K_\mathrm{total}$
Bellman, suboptimal at the start (because $\pi^{(1)}$ plans only $K_\mathrm{batch}$ steps
ahead even though $K_\mathrm{total} - K_\mathrm{batch}$ steps remain) and tight at the end.

\subsection{The value of a batched policy}

The expected reward under the batched policy starting from $\bb_0 = p_0$ is
\begin{equation}
V^\mathrm{batch}(p_0;\, \bc)
\;=\;
\E_{\bb_{K_\mathrm{batch}} \sim P(\,\cdot\,\mid p_0, \pi^{(1)},\, \bc)}\!
\bigl[\,V^\mathrm{batch}(\bb_{K_\mathrm{batch}};\, \bc, \text{remaining } n - 1 \text{ batches})\,\bigr],
\label{eq:Vbatch_recursive}
\end{equation}
unfolding to
\begin{equation}
V^\mathrm{batch}(p_0;\, \bc)
\;=\;
\sum_{\bb_{K_\mathrm{batch}}}\, P(\bb_{K_\mathrm{batch}} \mid p_0, \pi^{(1)},\, \bc)\,
V_0^{\mathrm{batch}_2}(\bb_{K_\mathrm{batch}};\, \bc),
\label{eq:Vbatch_2batch}
\end{equation}
in the two-batch case ($n = 2$), where $V_0^{\mathrm{batch}_2}$ denotes the standard
$K_\mathrm{batch}$ Bellman value at the new starting belief $\bb_{K_\mathrm{batch}}$ over
the second batch.  The sum is over reachable batch-boundary beliefs.

\subsection{The curse of enumeration over batch-boundary beliefs}

To compute Equation~\eqref{eq:Vbatch_2batch} exactly we would need
$V_0^{\mathrm{batch}_2}(\bb_{K_\mathrm{batch}};\, \bc)$ at every reachable
$\bb_{K_\mathrm{batch}}$.  In the count-tuple-keyed setting this is one fresh
$K_\mathrm{batch}$ Bellman solve per reachable count-tuple after batch~1.  At
$K_\mathrm{batch} = 4$ the reachable count-tuple set is $\sim 10^7$ in unfavourable
regimes, and each fresh solve costs $\sim 10$~seconds: enumeration cost is $10^8$ seconds
(thousands of CPU-years).

This is the same memo blow-up that capped exact $K_\mathrm{total} = K_\mathrm{batch} + 1$
in the first place, just relocated to the batch boundary.  Batched DP, naively
implemented, does not break the exact-DP ceiling.

\section{Stochastic batched DP}
\label{sec:stochastic}

The standard remedy for an intractable expectation over a large index set is to replace
exhaustive enumeration with Monte-Carlo sampling under the natural distribution
$P(\bb_{K_\mathrm{batch}} \mid p_0, \pi^{(1)},\, \bc)$.

\subsection{The estimator}

Sample $M$ trajectories under the current batched policy.  For each trajectory $m$ define
$\bb_{K_\mathrm{batch}}^{(m)}$ as the realized batch-1 terminal belief.  The estimator
\begin{equation}
\widehat{V}^\mathrm{batch}(p_0;\, \bc; \mathcal{M})
\;:=\;
\frac{1}{M} \sum_{m=1}^{M}
V_0^{\mathrm{batch}_2}\bigl(\bb_{K_\mathrm{batch}}^{(m)};\, \bc\bigr)
\label{eq:estimator}
\end{equation}
is an unbiased estimator of Equation~\eqref{eq:Vbatch_2batch}: the trajectories
$\mathcal{M} = \{\bb_{K_\mathrm{batch}}^{(m)}\}_{m=1}^M$ are i.i.d.\ samples from the
correct policy-induced distribution, and the integrand is computed exactly per sample by
solving fresh $K_\mathrm{batch}$ Bellman from $\bb_{K_\mathrm{batch}}^{(m)}$.

The variance of the estimator scales as $\Var[V_0^{\mathrm{batch}_2}] / M$ where the
variance is over the trajectory distribution; standard concentration inequalities give
the $1/\sqrt{M}$ confidence intervals.

\subsection{Per-sample gradients via the envelope theorem}

For outer optimization we need not just the value but its gradient with respect to $\bc$.
For each sampled trajectory we apply the sharp-max envelope theorem to the
$K_\mathrm{batch}$-step subproblem starting at $\bb_{K_\mathrm{batch}}^{(m)}$:
\begin{enumerate}[nosep, leftmargin=2em]
\item Run forward $K_\mathrm{batch}$ Bellman from $\bb_{K_\mathrm{batch}}^{(m)}$.  Record
      the optimal argmax action at every reachable belief in the second batch's subtree.
\item Re-evaluate $V_0^{\mathrm{batch}_2}(\bb_{K_\mathrm{batch}}^{(m)};\, \bc)$ by
      substituting the recorded argmaxes in place of the $\max$ operator.  This produces
      an argmax-free expression in $\bc$.
\item Apply reverse-mode automatic differentiation to the argmax-free expression.
\end{enumerate}

The envelope theorem certifies that the resulting gradient is exact:
\begin{equation}
\frac{\partial V_0^{\mathrm{batch}_2}(\bb_{K_\mathrm{batch}}^{(m)};\, \bc)}{\partial \bc}
\;=\;
\left.\frac{\partial}{\partial \bc}\right|_{\pi^\star_{batch_2} \text{ frozen}}
V_0^{\mathrm{batch}_2}(\bb_{K_\mathrm{batch}}^{(m)};\, \bc).
\label{eq:envelope}
\end{equation}
The first batch's contribution to the gradient (through how $\pi^{(1)}$ shapes the sampling
distribution of $\bb_{K_\mathrm{batch}}$) follows the same envelope-theorem-frozen recursion
applied to batch 1, and contributes a score-function term in $\partial \log P / \partial \bc$
that can be carried via reverse-mode AD on the trajectory's likelihood, exactly as in the
single-batch joint-DP gradient.

The aggregate gradient
\begin{equation}
\widehat{\nabla}_{\bc} V^\mathrm{batch}(p_0;\, \bc; \mathcal{M})
\;=\;
\frac{1}{M} \sum_{m=1}^{M}
\frac{\partial V^{(m)}(\bc)}{\partial \bc}
\label{eq:agg_grad}
\end{equation}
is unbiased for $\nabla_{\bc} V^\mathrm{batch}(p_0;\, \bc)$ and has the standard SGD
variance scaling.

\subsection{Outer SGD update}

The outer optimization becomes Adam (or any preferred stochastic optimizer) on $\bc$ with
mini-batches of size $M$:
\begin{equation}
\bc^{(t+1)} = \mathrm{Adam}\bigl(\bc^{(t)},\; \widehat{\nabla}_{\bc} V^\mathrm{batch}(p_0;\, \bc^{(t)}; \mathcal{M}^{(t)}),\; \eta\bigr).
\end{equation}
A fresh trajectory mini-batch $\mathcal{M}^{(t)}$ is sampled at each iteration $t$, so the
variance does not compound across iterations.  Standard SGD convergence theory applies:
under bounded gradient variance and standard step-size schedules the iterates converge in
expectation to a stationary point of $V^\mathrm{batch}(p_0;\, \bc)$.

\section{Gradient computation: a detailed derivation}
\label{sec:gradient}

This section derives the SBDP outer-hardware gradient explicitly.  We work the two-batch
case ($n = 2$) in full detail, identifying every $\bc$-dependence and unfolding each chain
rule.  We then generalize to $n > 2$ batches recursively.  No step is skipped; the goal is
that the formula in the algorithm box can be derived directly from first principles by
following this section.

\subsection{Setup for the two-batch case}
\label{sec:n2_setup}

For $n = 2$ with $K_\mathrm{total} = 2 K_\mathrm{batch}$, the batched value is
\begin{equation}
V^\mathrm{batch}(p_0;\, \bc)
\;=\;
\E_{\bx \sim p_0}\,
\E_{(\by_1, \ldots, \by_{K_\mathrm{batch}}) \mid \bx, \pi^{[1]}(\bc), \bc}
\!\left[\,
V^{[2]}\bigl(\bb_{K_\mathrm{batch}}(\bx, \by_{1:K_\mathrm{batch}}; \bc);\, \bc\bigr)
\,\right],
\label{eq:Vbatch_n2}
\end{equation}
where
\begin{itemize}[nosep,leftmargin=2em]
\item $\pi^{[1]}(\bc)$ is the batch-1 policy obtained by exact $K_\mathrm{batch}$ Bellman
      from $p_0$ at parameter $\bc$.  Each action in batch~1 is the local argmax
      $\bs_k = \pi^{[1]}(\bb_{k-1};\, \bc)$.
\item $\bb_{k} = T(\bb_{k-1}, \bs_k, \by_k;\, \bc)$ is the Bayesian belief update.
\item $V^{[2]}(\bb;\, \bc)$ is the standard $K_\mathrm{batch}$ Bellman value at starting
      belief $\bb$ with parameter $\bc$.  Because batch~2 is the last batch, its local
      terminal coincides with the global terminal $\Phi_0$.
\end{itemize}

The trajectory is the tuple $\tau := (\bx, \by_1, \ldots, \by_{K_\mathrm{batch}})$ with
joint density
\begin{equation}
P(\tau;\, \bc) \;=\; p_0(\bx) \cdot \prod_{k=1}^{K_\mathrm{batch}} p\bigl(\by_k \mid \bx, \bs_k(\tau_{<k}, \bc);\, \bc\bigr),
\label{eq:trajdensity}
\end{equation}
where $\bs_k(\tau_{<k}, \bc) = \pi^{[1]}(\bb_{k-1}(\tau_{<k}, \bc);\, \bc)$ is the
deterministic action under the batch-1 policy.

We want $\nabla_{\bc}\, V^\mathrm{batch}(p_0;\, \bc)$.

\subsection{Identifying every \texorpdfstring{$\bc$}{c}-dependence}
\label{sec:c_dependences}

Equation~\eqref{eq:Vbatch_n2} has $\bc$ entering through five distinct paths.  We list
them explicitly because each contributes to the gradient differently:

\begin{enumerate}[nosep,leftmargin=2em]
\item[(P1)] \emph{Batch-1 policy actions} $\bs_k(\tau_{<k}, \bc)$ via the argmax of the local
      $K_\mathrm{batch}$ Bellman.
\item[(P2)] \emph{Within-batch-1 observation likelihood} $p(\by_k \mid \bx, \bs_k;\, \bc)$.
\item[(P3)] \emph{Within-batch-1 belief update} $T(\bb_{k-1}, \bs_k, \by_k;\, \bc)$, which
      shapes how $\bb_{K_\mathrm{batch}}$ depends on the trajectory.
\item[(P4)] \emph{Batch-2 Bellman value at fixed starting belief}
      $V^{[2]}(\cdot;\, \bc)$, which depends on $\bc$ through batch-2's likelihood and
      belief updates.
\item[(P5)] \emph{Starting belief input to batch-2}, $\bb_{K_\mathrm{batch}}(\bx, \by_{1:K_\mathrm{batch}}; \bc)$.
      Evaluating $V^{[2]}$ at a starting belief that itself depends on $\bc$ pulls a chain
      of (P3) sensitivities in.
\end{enumerate}

We now show that (P1) contributes zero to the gradient under hard argmax, while
(P2)--(P5) contribute through two channels: a pathwise term (carrying P3, P4, P5) and a
score-function term (carrying P2).

\paragraph{Hard-argmax policy contributes zero.}
The batch-1 policy is $\bs_k(\bb_{k-1}; \bc) = \argmax_{\bs} Q^{[1]}_k(\bb_{k-1}, \bs;\, \bc)$.
For an interior strict maximizer, the argmax is locally constant in $\bc$:
$\partial \bs_k / \partial \bc = 0$ almost everywhere in $\bc$.  Hence path (P1)
contributes a measure-zero set of jump points, which integrate to zero against
$\bc$-perturbations.  Throughout, we therefore treat the realized actions
$\bs^{(m)}_{1:K_\mathrm{batch}}$ on each sampled trajectory as
$\bc$-frozen at their recorded values.  This is the same locally-constant-argmax
assumption used by the sharp-max envelope theorem of the joint-DP framework, applied here
to the policy choices.

\subsection{The REINFORCE / pathwise decomposition}
\label{sec:reinforce}

Let $F(\tau;\, \bc) := V^{[2]}(\bb_{K_\mathrm{batch}}(\tau, \bc);\, \bc)$.  Then
$V^\mathrm{batch}(p_0;\, \bc) = \E_{\tau \sim P(\cdot;\, \bc)}[F(\tau;\, \bc)]$.  By the
log-derivative identity ($\nabla_{\bc} P = P \cdot \nabla_{\bc} \log P$),
\begin{equation}
\nabla_{\bc}\, \E_{\tau \sim P(\cdot;\, \bc)}\!\bigl[F(\tau;\, \bc)\bigr]
\;=\;
\E_{\tau \sim P(\cdot;\, \bc)}\!\bigl[\,\nabla_{\bc} F(\tau;\, \bc)\,\bigr]
\;+\;
\E_{\tau \sim P(\cdot;\, \bc)}\!\bigl[\,F(\tau;\, \bc)\, \nabla_{\bc} \log P(\tau;\, \bc)\,\bigr].
\label{eq:reinforce}
\end{equation}
The first term is the \emph{pathwise gradient} (deterministic chain rule on the integrand
at fixed sampled $\tau$).  The second term is the \emph{score-function gradient}
(REINFORCE-style sensitivity of the trajectory distribution).  Both are needed for an
unbiased estimator.

\subsection{Pathwise gradient \texorpdfstring{$\nabla_{\bc} F(\tau;\, \bc)$}{nabla\_c F}}
\label{sec:pathwise}

By chain rule on $F = V^{[2]}\bigl(\bb_{K_\mathrm{batch}}(\tau, \bc);\, \bc\bigr)$,
\begin{equation}
\nabla_{\bc} F(\tau;\, \bc)
\;=\;
\underbrace{\Bigl[\,\nabla_{\bb} V^{[2]}\!\bigl(\bb_{K_\mathrm{batch}};\, \bc\bigr)\,\Bigr]^\top \cdot \nabla_{\bc} \bb_{K_\mathrm{batch}}(\tau, \bc)}_{\text{(A): input-belief sensitivity}}
\;+\;
\underbrace{\biggl.\frac{\partial V^{[2]}\bigl(\bb;\, \bc\bigr)}{\partial \bc}\biggr|_{\bb \text{ held fixed at } \bb_{K_\mathrm{batch}}}}_{\text{(B): standard joint-DP gradient}}.
\label{eq:pathwise}
\end{equation}
Each term has a concrete computation; we unfold them in turn.

\paragraph{Term (B): standard joint-DP gradient at fixed starting belief.}

This is the existing joint-DP gradient of the paper's main framework, applied to
$V^{[2]}$.  Specifically, holding the starting belief fixed at $\bb_{K_\mathrm{batch}}^{(m)}$
on the sampled trajectory, we run forward the $K_\mathrm{batch}$ Bellman recursion at
parameter $\bc$ to obtain the policy memo for batch~2.  Then we apply the sharp-max
envelope theorem: freeze the recorded action argmaxes throughout the batch-2 subtree, and
compute reverse-mode automatic differentiation on the resulting argmax-free expression for
$V^{[2]}(\bb_{K_\mathrm{batch}}^{(m)};\, \bc)$.  The result is exact for the batch-2
subproblem.

In the count-tuple-keyed setting, term (B) is computable in time $O(N_\mathrm{batch-2-memo} \cdot \mathrm{cost}(p, T))$, comparable to one $K_\mathrm{batch}$ Bellman solve.

\paragraph{Term (A): input-belief sensitivity through belief-update chain.}

This is two distinct factors:
\begin{itemize}[nosep, leftmargin=2em]
\item $\nabla_{\bb} V^{[2]}(\bb;\, \bc)$: the gradient of the batch-2 Bellman value with
      respect to the \emph{input belief} $\bb$, evaluated at $\bb = \bb_{K_\mathrm{batch}}^{(m)}$.
\item $\nabla_{\bc} \bb_{K_\mathrm{batch}}(\tau, \bc)$: the sensitivity of the batch-1
      terminal belief to $\bc$ along the realized batch-1 trajectory.
\end{itemize}

\textbf{Computing $\nabla_{\bb} V^{[2]}$.}  The batch-2 Bellman recursion at starting
belief $\bb$ contains $\bb$ in two places: (i) directly, as the argument of the value
function; (ii) through the marginal observation probability $\Pr(\by \mid \bb, \bs;\, \bc) =
\sum_{\bx} \bb(\bx)\, p(\by \mid \bx, \bs;\, \bc)$, which appears in every Bellman backup
weighting future values.  In the count-tuple-keyed framework with $\bb$ on a $K_\Phi$-grid,
$\bb$ is a $K_\Phi$-dimensional vector; $\nabla_{\bb} V^{[2]}$ is a $K_\Phi$-vector.  It is
computable by reverse-mode AD through the argmax-frozen Bellman recursion, treating $\bb$
as the differentiation target instead of $\bc$.

\textbf{Computing $\nabla_{\bc} \bb_{K_\mathrm{batch}}$.}  The belief at the end of batch~1
is the recursion
\[
\bb_k = T(\bb_{k-1}, \bs_k, \by_k;\, \bc), \qquad k = 1, \ldots, K_\mathrm{batch},
\]
with $\bs_k$ and $\by_k$ frozen at their recorded values.  Differentiating w.r.t. $\bc$
with frozen $(\bs_k, \by_k)$:
\begin{equation}
\frac{\partial \bb_k}{\partial \bc}
\;=\;
\biggl.\frac{\partial T}{\partial \bc}\biggr|_{\bb_{k-1}, \bs_k, \by_k}
\;+\;
\biggl.\frac{\partial T}{\partial \bb}\biggr|_{\bb_{k-1}, \bs_k, \by_k}
\cdot
\frac{\partial \bb_{k-1}}{\partial \bc},
\label{eq:belief_chain}
\end{equation}
unfolding from $\partial \bb_0 / \partial \bc = 0$ (the prior is $\bc$-independent).
This is a standard chain-rule recursion through the belief-update operator and is
computable by one reverse-mode AD pass over the $K_\mathrm{batch}$-step chain.

\textbf{Combining (A1) and (A2).}  Term~(A) of Equation~\eqref{eq:pathwise} is
\[
\sum_{i=1}^{K_\Phi} \frac{\partial V^{[2]}}{\partial \bb_{K_\mathrm{batch}, i}}\bigl(\bb_{K_\mathrm{batch}}; \bc\bigr) \cdot \frac{\partial \bb_{K_\mathrm{batch}, i}}{\partial \bc}.
\]
Both factors are vectors in $\R^{K_\Phi}$; the inner product collapses them into a scalar
gradient w.r.t. $\bc$.  In practice the two reverse-mode AD passes are combined into one,
treating the entire pipeline (batch-1 belief updates feeding batch-2 Bellman value) as a
single computational graph with $\bc$ as the differentiation target.

\subsection{Score-function gradient \texorpdfstring{$\nabla_{\bc} \log P(\tau;\, \bc)$}{nabla\_c log P}}
\label{sec:score_function}

From Equation~\eqref{eq:trajdensity},
\[
\log P(\tau;\, \bc)
\;=\;
\log p_0(\bx)
\;+\;
\sum_{k=1}^{K_\mathrm{batch}} \log p\bigl(\by_k \mid \bx, \bs_k;\, \bc\bigr).
\]
The prior $p_0(\bx)$ is $\bc$-independent.  Differentiating:
\begin{equation}
\nabla_{\bc} \log P(\tau;\, \bc)
\;=\;
\sum_{k=1}^{K_\mathrm{batch}} \nabla_{\bc} \log p\bigl(\by_k \mid \bx, \bs_k;\, \bc\bigr),
\label{eq:scorefn}
\end{equation}
with $\bs_k$ frozen by the hard-argmax property of \S\ref{sec:c_dependences}.  Each summand
is the per-shot likelihood score, computable by reverse-mode AD on the closed-form
likelihood $p(\by \mid \bx, \bs;\, \bc)$.  Note that this is a vector in $\R^{d_c}$ (one
component per hardware coordinate); the score-function gradient is a sum of such vectors
across batch-1 shots.

\subsection{The full per-trajectory gradient estimator}
\label{sec:full_estimator}

Combining \eqref{eq:reinforce}, \eqref{eq:pathwise}, and \eqref{eq:scorefn}, the
per-trajectory gradient estimator is
\begin{equation}
g^{(m)}
\;=\;
\underbrace{\Bigl[\nabla_{\bb} V^{[2]}\bigr|_{\bb_{K_\mathrm{batch}}^{(m)}}\!\Bigr]^\top \cdot \nabla_{\bc} \bb_{K_\mathrm{batch}}^{(m)}}_{\text{Term A: belief-input × belief chain}}
\;+\;
\underbrace{\biggl.\frac{\partial V^{[2]}\bigl(\bb;\, \bc\bigr)}{\partial \bc}\biggr|_{\bb = \bb_{K_\mathrm{batch}}^{(m)}}}_{\text{Term B: joint-DP gradient at fixed input}}
\;+\;
\underbrace{V^{[2]}\bigl(\bb_{K_\mathrm{batch}}^{(m)};\, \bc\bigr) \cdot \sum_{k=1}^{K_\mathrm{batch}} \nabla_{\bc} \log p\bigl(\by_k^{(m)} \mid \bx^{(m)}, \bs_k^{(m)};\, \bc\bigr)}_{\text{Term C: score-function correction}}.
\label{eq:gm}
\end{equation}
Aggregating across $M$ samples,
\begin{equation}
\widehat{\nabla}_{\bc}\, V^\mathrm{batch}(p_0;\, \bc; \mathcal{M})
\;=\;
\frac{1}{M} \sum_{m=1}^{M} g^{(m)}.
\label{eq:full_estimator}
\end{equation}
This estimator is unbiased for $\nabla_{\bc}\, V^\mathrm{batch}(p_0;\, \bc)$ by direct
substitution of \eqref{eq:reinforce} and the law of large numbers.

\paragraph{Roles of the three terms.}  Term~B is what one would write down naively if the
trajectory were $\bc$-independent: the standard joint-DP gradient evaluated at the sampled
batch-2 starting belief.  Term~A corrects for the fact that $\bb_{K_\mathrm{batch}}$ itself
depends on $\bc$ through the chain of belief updates that built it.  Term~C corrects for
the fact that the trajectory's sampling distribution depends on $\bc$ through the
likelihood; without it the estimator is biased toward the trajectories that current $\bc$
makes likely.

\paragraph{Failure modes if any term is dropped.}  Dropping Term~A produces a bias
proportional to the rate at which $\bb_{K_\mathrm{batch}}$ shifts under $\bc$, scaled by
the slope of $V^{[2]}$ in $\bb$.  Dropping Term~C produces a bias proportional to the
score-function-times-value covariance.  Both are problem-dependent; in cases where the
likelihood depends on $\bc$ only weakly across the prior support, Terms~A and~C can be
small, but they are not zero in general.

\subsection{Variance reduction with baselines}
\label{sec:baseline}

The score-function term~C in \eqref{eq:gm} is the dominant variance contributor, because
it is a product of two random quantities (the value $V^{[2]}$ and the per-shot scores).
Standard variance reduction subtracts a $\bc$-only baseline $B(\bc)$ from $V^{[2]}$ before
multiplying by the score:
\begin{equation}
g^{(m)}_{\mathrm{reduced}}
\;=\;
\text{(A)} + \text{(B)} + \bigl[V^{[2]}(\bb_{K_\mathrm{batch}}^{(m)};\, \bc) - B(\bc)\bigr] \cdot \sum_k \nabla_{\bc} \log p(\by_k^{(m)} \mid \bx^{(m)}, \bs_k^{(m)}; \bc).
\label{eq:gm_baseline}
\end{equation}
The estimator remains unbiased for any $B(\bc)$ that does not depend on $\tau^{(m)}$,
because $\E_{\tau}[\nabla_{\bc} \log P(\tau;\, \bc)] = 0$ identically.

\paragraph{Choice of baseline.}  Two standard choices:
\begin{itemize}[nosep, leftmargin=2em]
\item \emph{Empirical mean:} $B(\bc) = \frac{1}{M} \sum_{m'=1}^M V^{[2]}(\bb_{K_\mathrm{batch}}^{(m')};\, \bc)$.
      Cheap (free, given the values are already computed) but biased if used naively
      (because $B$ depends on the same samples whose scores it modifies).  The standard
      remedy is leave-one-out: $B^{(m)}(\bc) = \frac{1}{M-1} \sum_{m' \ne m} V^{[2]}(\bb_{K_\mathrm{batch}}^{(m')};\, \bc)$.
\item \emph{Local Bellman value:} $B(\bc) = V^{[1]}_{\mathrm{local}}(p_0;\, \bc)$, the
      $K_\mathrm{batch}$ Bellman value of batch~1 at $p_0$ with local terminal $\Phi_0$.
      This is computable for free during the batch-1 Bellman solve and is the optimal
      baseline in the sense of minimum variance among $\bc$-only constants when the local
      terminal of batch~1 matches the global terminal of $V^{[2]}$.
\end{itemize}

The local-Bellman baseline is the natural ``advantage-function'' choice and we recommend
it when the local and global terminals coincide (e.g., $\Phi_0 = -\mathrm{Var}(\bx \mid \bb)$
applied at end of every batch including the final one).

\subsection{Generalization to \texorpdfstring{$n > 2$}{n > 2} batches}
\label{sec:nbatch_general}

We now extend the derivation of \S\ref{sec:n2_setup}--\ref{sec:full_estimator} to general
$n$ batches.  The structure mirrors the two-batch case exactly; the additional content is
the per-batch decomposition of the score-function term and the corresponding advantage-form
variance reduction.

\subsubsection{Trajectory factorization across batches}
\label{sec:nbatch_factorization}

For $n$ batches with $K_\mathrm{total} = n K_\mathrm{batch}$, denote the per-batch
sub-trajectories
\[
\tau_j := \bigl(\by_{(j-1)K_\mathrm{batch} + 1}, \ldots, \by_{j K_\mathrm{batch}}\bigr), \qquad j = 1, \ldots, n,
\]
and the full trajectory $\tau := (\bx, \tau_1, \ldots, \tau_n)$.  Let
$\bb^{\mathrm{bdy}}_j := \bb_{j K_\mathrm{batch}}$ denote the belief at the end of batch
$j$ (so $\bb^{\mathrm{bdy}}_0 = p_0$ and $\bb^{\mathrm{bdy}}_n = \bb_{K_\mathrm{total}}$).

The trajectory density factorizes as a single product over all $K_\mathrm{total}$ shots,
\begin{equation}
P(\tau;\, \bc)
\;=\;
p_0(\bx) \cdot \prod_{k=1}^{K_\mathrm{total}} p\bigl(\by_k \mid \bx, \bs_k;\, \bc\bigr),
\label{eq:Pfactor_full}
\end{equation}
which by grouping shots into batches becomes
\begin{equation}
P(\tau;\, \bc)
\;=\;
p_0(\bx) \cdot \prod_{j=1}^{n} P_j\bigl(\tau_j \mid \bx, \bb^{\mathrm{bdy}}_{j-1};\, \bc\bigr),
\quad \text{where} \quad
P_j(\tau_j \mid \bx, \bb^{\mathrm{bdy}}_{j-1};\, \bc) = \prod_{k \in \text{batch } j} p(\by_k \mid \bx, \bs_k;\, \bc).
\label{eq:Pfactor_perbatch}
\end{equation}
The per-batch density $P_j$ depends on $\bc$ through both the within-batch likelihoods and
the local Bellman policy that selects $\bs_k$ for $k \in$~batch~$j$.  As before, the policy
contribution vanishes under hard argmax, so $\bs_k$ values are frozen at recorded argmaxes.

\subsubsection{REINFORCE applied to the full trajectory}
\label{sec:nbatch_reinforce}

Because the full trajectory factorization \eqref{eq:Pfactor_full} contains all $K_\mathrm{total}$
shots in one product, the REINFORCE identity \eqref{eq:reinforce} applies to the full
trajectory directly.  With $F(\tau;\, \bc) := \Phi_0(\bb_{K_\mathrm{total}}(\tau, \bc);\, \bc)$,
\begin{equation}
\nabla_{\bc} V^\mathrm{batch}(p_0;\, \bc)
\;=\;
\E_{\tau \sim P(\cdot;\, \bc)}\!\bigl[\,\nabla_{\bc} F(\tau;\, \bc)\,\bigr]
\;+\;
\E_{\tau \sim P(\cdot;\, \bc)}\!\bigl[\,F(\tau;\, \bc) \cdot \nabla_{\bc} \log P(\tau;\, \bc)\,\bigr].
\label{eq:nbatch_reinforce}
\end{equation}
Both terms unfold across all $K_\mathrm{total}$ steps; the per-batch decomposition arises
naturally because the underlying recursion has batch boundaries built into the policy.

\subsubsection{Pathwise gradient across all \texorpdfstring{$K_\mathrm{total}$}{K\_total} steps}
\label{sec:nbatch_pathwise}

The terminal reward $\Phi_0(\bb_{K_\mathrm{total}}(\tau, \bc);\, \bc)$ depends on $\bc$
through (i) any direct $\bc$-dependence in the functional form of $\Phi_0$ and (ii) the
chain of $K_\mathrm{total}$ belief updates.  By chain rule,
\begin{equation}
\nabla_{\bc} F(\tau;\, \bc)
\;=\;
\biggl.\frac{\partial \Phi_0(\bb;\, \bc)}{\partial \bc}\biggr|_{\bb = \bb_{K_\mathrm{total}}}
\;+\;
\Bigl[\nabla_{\bb} \Phi_0(\bb_{K_\mathrm{total}};\, \bc)\Bigr]^\top \cdot \nabla_{\bc}\, \bb_{K_\mathrm{total}}(\tau, \bc),
\label{eq:nbatch_pathwise}
\end{equation}
with the belief-chain derivative
\begin{equation}
\frac{\partial \bb_k}{\partial \bc}
\;=\;
\biggl.\frac{\partial T}{\partial \bc}\biggr|_{\bb_{k-1}, \bs_k, \by_k}
\;+\;
\biggl.\frac{\partial T}{\partial \bb}\biggr|_{\bb_{k-1}, \bs_k, \by_k}
\cdot
\frac{\partial \bb_{k-1}}{\partial \bc},
\qquad k = 1, \ldots, K_\mathrm{total},
\label{eq:nbatch_belief_chain}
\end{equation}
unfolding from $\partial \bb_0 / \partial \bc = 0$.  Equation~\eqref{eq:nbatch_belief_chain}
is the same belief-chain recursion as Equation~\eqref{eq:belief_chain} but extended over all
$K_\mathrm{total}$ steps, not just the first $K_\mathrm{batch}$.  In implementation it is
one reverse-mode AD pass over the full trajectory's belief-update chain.  Note that
the batch boundaries do not appear specially in this term: the belief update is the same
operator at every step, regardless of which batch the step belongs to.

\paragraph{Comparison with the two-batch pathwise term.}  In the two-batch derivation
\eqref{eq:pathwise} the pathwise term split into an "input belief" piece~(A) and a
"fixed-input batch-2 Bellman gradient" piece~(B).  Here we have folded both pieces into
a single chain rule expression \eqref{eq:nbatch_pathwise}.  The split is equivalent: the
"fixed-input batch-2 Bellman gradient" of the two-batch case corresponds to differentiating
the second batch's belief-update chain with the second batch's input frozen, while the
"input belief sensitivity" corresponds to differentiating the first batch's chain.  Both
flatten into the same total chain over $K_\mathrm{total}$ steps.

\subsubsection{Per-batch decomposition of the score-function term}
\label{sec:nbatch_score}

The full score function unfolds across all shots:
\begin{equation}
\nabla_{\bc} \log P(\tau;\, \bc)
\;=\;
\sum_{k=1}^{K_\mathrm{total}} \nabla_{\bc} \log p\bigl(\by_k \mid \bx, \bs_k;\, \bc\bigr)
\;=\;
\sum_{j=1}^{n} \underbrace{\nabla_{\bc} \log P_j\bigl(\tau_j \mid \bx, \bb^{\mathrm{bdy}}_{j-1};\, \bc\bigr)}_{\sum_{k \in \text{batch } j} \nabla_{\bc} \log p(\by_k \mid \bx, \bs_k;\, \bc)}.
\label{eq:nbatch_score_decomp}
\end{equation}
The grouping into $n$ per-batch contributions is purely notational at this stage — the
per-batch sums are just additive.  But it becomes load-bearing in the next step, where
each per-batch term gets multiplied by the residual remaining-horizon value from that
batch onwards rather than by the global terminal.

\subsubsection{Equivalent per-batch score-function form (for variance reduction)}
\label{sec:nbatch_perbatch}

The "raw" REINFORCE gradient \eqref{eq:nbatch_reinforce} multiplies the score function by
the global terminal $F(\tau;\, \bc) = \Phi_0(\bb_{K_\mathrm{total}};\, \bc)$.  By the same
identity that proves $\E_\tau[\nabla_{\bc} \log P(\tau;\, \bc)] = 0$ for an arbitrary
constant baseline, we can substitute any \emph{$\bc$-only} or \emph{causally-prior}
quantity for the score-function multiplier within each batch's term, without introducing
bias.  Specifically:
\begin{lemma}[Per-batch score-function multiplier substitution]
\label{lem:perbatch_substitution}
For each $j = 1, \ldots, n$, let $W_j(\bc, \tau_{<j})$ be any function depending on $\bc$
and on the trajectory \emph{prior} to batch $j$ (i.e., on $\bx$ and on
$\tau_1, \ldots, \tau_{j-1}$).  Then
\begin{equation}
\E_{\tau} \bigl[\,W_j(\bc, \tau_{<j}) \cdot \nabla_{\bc} \log P_j(\tau_j \mid \bx, \bb^{\mathrm{bdy}}_{j-1};\, \bc)\,\bigr] \;=\; 0.
\end{equation}
\end{lemma}

\noindent\emph{Proof.}  By the tower property of conditional expectation,
\[
\E_{\tau}\!\bigl[W_j \cdot \nabla_{\bc} \log P_j\bigr]
\;=\;
\E_{\tau_{<j}}\!\Bigl[\,W_j(\bc, \tau_{<j}) \cdot \E_{\tau_j \mid \tau_{<j}}\!\bigl[\nabla_{\bc} \log P_j(\tau_j \mid \bx, \bb^{\mathrm{bdy}}_{j-1};\, \bc)\bigr]\,\Bigr],
\]
and the inner conditional expectation is zero because $\E[\nabla_{\bc} \log p] = 0$ for any
parametric density.  $\square$

\medskip

This means we can replace the score-function multiplier inside batch $j$'s term with the
\emph{remaining-horizon value} $V_{\mathrm{global}}^{(j+1)}(\bb^{\mathrm{bdy}}_j;\, \bc)$
plus or minus any $\bc$-only quantity:
\begin{equation}
\E_\tau\bigl[F \cdot \nabla_{\bc} \log P_j\bigr]
\;=\;
\E_\tau\bigl[V_{\mathrm{global}}^{(j+1)}(\bb^{\mathrm{bdy}}_j;\, \bc) \cdot \nabla_{\bc} \log P_j\bigr],
\label{eq:perbatch_value_substitution}
\end{equation}
because the difference $F - V_{\mathrm{global}}^{(j+1)}(\bb^{\mathrm{bdy}}_j;\, \bc)$ has
zero conditional expectation given $\tau_{<j+1}$ (it is exactly the residual stochasticity
in batches $j+1, \ldots, n$ given the boundary belief), and that residual is independent of
batch $j$'s score function only when integrated against batch $j$'s score (which is what
Lemma~\ref{lem:perbatch_substitution} certifies).

Combining \eqref{eq:nbatch_reinforce}, \eqref{eq:nbatch_score_decomp}, and
\eqref{eq:perbatch_value_substitution} yields the \emph{per-batch} REINFORCE form:
\begin{equation}
\boxed{\;
\nabla_{\bc} V^\mathrm{batch}(p_0;\, \bc)
\;=\;
\E_\tau\!\biggl[\,
\nabla_{\bc} F(\tau;\, \bc)
\;+\;
\sum_{j=1}^{n} V_{\mathrm{global}}^{(j+1)}\!\bigl(\bb^{\mathrm{bdy}}_j;\, \bc\bigr) \cdot \nabla_{\bc} \log P_j(\tau_j \mid \bx, \bb^{\mathrm{bdy}}_{j-1};\, \bc)
\,\biggr].
\;}
\label{eq:nbatch_main}
\end{equation}
Here $V_{\mathrm{global}}^{(n+1)} := \Phi_0(\bb;\, \bc)$ by convention (so the $j = n$ term
collapses to the standard score function multiplied by the realized terminal).

\subsubsection{Variance reduction with per-batch advantage baselines}
\label{sec:nbatch_advantage}

By the same lemma, we can subtract any $\bc$-only or causally-prior quantity from
$V_{\mathrm{global}}^{(j+1)}(\bb^{\mathrm{bdy}}_j;\, \bc)$ inside the $j$-th term without
introducing bias.  Define the per-batch advantage
\begin{equation}
A_j(\bb^{\mathrm{bdy}}_j;\, \bc) \;:=\; V_{\mathrm{global}}^{(j+1)}\!\bigl(\bb^{\mathrm{bdy}}_j;\, \bc\bigr) - V_{\mathrm{global}}^{(j)}\!\bigl(\bb^{\mathrm{bdy}}_{j-1};\, \bc\bigr).
\label{eq:advantage}
\end{equation}
This is the \emph{realized-minus-expected} value of batch $j$: it measures how much
unexpectedly more (or less) value batch $j$ delivered relative to what was anticipated
from $\bb^{\mathrm{bdy}}_{j-1}$ at the start of the batch.  By
Lemma~\ref{lem:perbatch_substitution}, $\E[V_{\mathrm{global}}^{(j)}(\bb^{\mathrm{bdy}}_{j-1};\, \bc) \cdot \nabla_{\bc} \log P_j] = 0$, so substituting $V_{\mathrm{global}}^{(j+1)}$ with $A_j$ leaves the gradient unbiased:
\begin{equation}
\nabla_{\bc} V^\mathrm{batch}(p_0;\, \bc)
\;=\;
\E_\tau\!\Bigl[\,
\nabla_{\bc} F(\tau;\, \bc)
\;+\;
\sum_{j=1}^{n} A_j(\bb^{\mathrm{bdy}}_j;\, \bc) \cdot \nabla_{\bc} \log P_j(\tau_j \mid \bx, \bb^{\mathrm{bdy}}_{j-1};\, \bc)
\,\Bigr].
\label{eq:nbatch_advantage}
\end{equation}
Equation~\eqref{eq:nbatch_advantage} is the actor-critic / advantage form.  Variance is
typically much smaller than \eqref{eq:nbatch_main} because $A_j$ has zero conditional mean
(by definition), while $V_{\mathrm{global}}^{(j+1)}$ has the full per-trajectory variance.

\paragraph{Estimating $V_{\mathrm{global}}^{(j)}$.}  The advantage form requires
$V_{\mathrm{global}}^{(j)}$ at the boundary beliefs, but we do not have a closed form for
these in the general SBDP setting.  Three practical estimators:
\begin{itemize}[nosep, leftmargin=2em]
\item \emph{Empirical mean over samples}: $\widehat{V}_{\mathrm{global}}^{(j)}(\bb;\, \bc) \approx
  \frac{1}{M} \sum_{m'=1}^{M} F(\tau^{(m')};\, \bc)$, treating samples as i.i.d. from the trajectory
  distribution.  Cheap; introduces a small leave-one-out bias unless one excludes the current
  sample's contribution.
\item \emph{Local Bellman value}: $\widehat{V}_{\mathrm{global}}^{(j)}(\bb;\, \bc) \approx
  V^{[j]}_\mathrm{local}(\bb;\, \bc)$, the standard $K_\mathrm{batch}$ Bellman value with terminal
  $\Phi_0$ at end of batch $j$.  Free during the local Bellman solve.  Optimal as a baseline
  when local terminal coincides with global terminal.
\item \emph{Bootstrap from sampled trajectories}: simulate one further batch under the current
  policy starting from $\bb^{\mathrm{bdy}}_j$ and record its terminal value.  Equivalent to
  one-step temporal-difference estimation.
\end{itemize}
Any of these is unbiased as a baseline (Lemma~\ref{lem:perbatch_substitution}) but they trade
off variance vs.\ implementation cost.

\subsubsection{Per-trajectory gradient estimator for general \texorpdfstring{$n$}{n}}
\label{sec:nbatch_estimator}

Combining \eqref{eq:nbatch_pathwise} and \eqref{eq:nbatch_advantage}, the per-trajectory
gradient estimator with advantage variance reduction is
\begin{equation}
g^{(m)}_n
\;=\;
\underbrace{\biggl[\frac{\partial \Phi_0}{\partial \bc} + (\nabla_{\bb} \Phi_0)^\top \cdot \nabla_{\bc} \bb_{K_\mathrm{total}}^{(m)}\biggr]}_{\text{pathwise (single chain over } K_\mathrm{total}\text{)}}
\;+\;
\sum_{j=1}^{n} \underbrace{A_j^{(m)} \cdot \sum_{k \in \text{batch } j} \nabla_{\bc} \log p\bigl(\by_k^{(m)} \mid \bx^{(m)}, \bs_k^{(m)};\, \bc\bigr)}_{\text{per-batch advantage} \times \text{batch-} j \text{ score}}.
\label{eq:gm_nbatch_full}
\end{equation}
Aggregating across $M$ samples:
\begin{equation}
\widehat{\nabla}_{\bc}\, V^\mathrm{batch}(p_0;\, \bc; \mathcal{M})
\;=\;
\frac{1}{M} \sum_{m=1}^{M} g^{(m)}_n.
\end{equation}
This is unbiased for $\nabla_{\bc} V^\mathrm{batch}$ for any choice of advantage baselines
$\widehat{V}_{\mathrm{global}}^{(j)}$ that depend only on $\bc$ and $\tau_{<j}^{(m)}$.

\paragraph{Per-step cost decomposition.}  Per sample $m$:
\begin{itemize}[nosep, leftmargin=2em]
\item One reverse-mode AD pass over the full $K_\mathrm{total}$-step belief-update chain:
      pathwise term.  Cost: $O(K_\mathrm{total} \cdot K_\Phi \cdot d_c)$.
\item For each of $n$ batch boundaries: one $K_\mathrm{batch}$ Bellman solve to obtain
      $V^{[j]}_\mathrm{local}$ for the local-Bellman baseline (or alternative).  Cost:
      $O(n \cdot \mathrm{cost}(K_\mathrm{batch} \text{ Bellman}))$.
\item One additional reverse-mode AD pass for each of the $K_\mathrm{total}$ score-function
      terms: cost $O(K_\mathrm{total} \cdot d_c)$.
\end{itemize}
Total per sample: dominated by $n \cdot \mathrm{cost}(K_\mathrm{batch} \text{ Bellman})$ for
typical problems.  Aggregate over $M$ samples in parallel.

\subsubsection{Sanity check: \texorpdfstring{$n = 2$}{n=2} reduces to \S\ref{sec:full_estimator}}

At $n = 2$, the per-batch sum in \eqref{eq:gm_nbatch_full} contains two terms:
$j = 1$ (advantage of batch~1) and $j = 2$ (advantage of batch~2).
\begin{itemize}[nosep, leftmargin=2em]
\item $A_1 = V_{\mathrm{global}}^{(2)}(\bb^{\mathrm{bdy}}_1;\, \bc) - V_{\mathrm{global}}^{(1)}(p_0;\, \bc)$.
      The second term is $\bc$-only and frequently dropped; the first term is exactly the
      $V^{[2]}(\bb_{K_\mathrm{batch}}^{(m)};\, \bc)$ that appeared as Term~C of
      Equation~\eqref{eq:gm}.
\item $A_2 = V_{\mathrm{global}}^{(3)}(\bb^{\mathrm{bdy}}_2;\, \bc) - V_{\mathrm{global}}^{(2)}(\bb^{\mathrm{bdy}}_1;\, \bc) = \Phi_0(\bb_{K_\mathrm{total}};\, \bc) - V^{[2]}(\bb^{\mathrm{bdy}}_1;\, \bc)$.
      The first piece is the realized terminal; the second is the predicted-from-batch-1-end value.
\end{itemize}
The pathwise gradient \eqref{eq:nbatch_pathwise} over the $K_\mathrm{total} = 2 K_\mathrm{batch}$
chain coincides with Term~A + Term~B of Equation~\eqref{eq:pathwise}.  The score-function term
multiplied by $A_1$ corresponds to the score-function-times-baseline-corrected version of
Term~C.  The score-function term multiplied by $A_2$ contains the batch-2 score-function
contribution (which was implicitly handled inside Term~B in the two-batch derivation, since
Term~B used exact within-batch Bellman to integrate over batch-2's outcomes analytically).

In short: the $n = 2$ case of \eqref{eq:gm_nbatch_full} is equivalent to
Equation~\eqref{eq:gm} after baseline shifting and after recognizing that the batch-2
score-function contribution is implicitly captured inside the batch-2 exact Bellman gradient.
The two derivations agree.

\paragraph{Why one might prefer $n = 2$'s explicit form.}  In the $n = 2$ derivation we kept
the second batch's exact-Bellman-with-summation-over-outcomes structure as Term~B; this
avoids one of the two stochastic score-function contributions and reduces variance.  For
general $n$, the analog is to use exact Bellman within the \emph{last} batch (which by
definition has terminal aligned with the global terminal) and stochastic batched DP only
for batch boundaries $j = 1, \ldots, n-1$.  This hybrid keeps batch~$n$'s contribution
deterministic and fully Bellman-exact, identical to the existing joint-DP gradient
machinery, while only sampling at boundaries $j < n$.

\subsection{Connection back to single-batch joint-DP}
\label{sec:single_batch_check}

For $n = 1$ (single batch with $K_\mathrm{total} = K_\mathrm{batch}$), there is no
batch-boundary sampling: the entire integral over $\tau$ is performed exactly inside the
Bellman recursion via the marginal observation probability sums.
Equation~\eqref{eq:gm_n} reduces to
\[
\nabla_{\bc} V^\mathrm{batch}(p_0;\, \bc)
\;=\;
\nabla_{\bc} V^{[1]}(p_0;\, \bc)
\;=\;
\biggl.\frac{\partial V^{[1]}(\bb;\, \bc)}{\partial \bc}\biggr|_{\bb = p_0, \pi^\star \text{ frozen}},
\]
which is exactly the sharp-max envelope-theorem gradient of the paper's main framework.
The score-function term vanishes because there are no batch boundaries and the $\E_\tau$
integration is performed exactly.  This recovers the exact-Bellman joint-DP gradient as
the $n = 1$ specialization of SBDP.  For $n \ge 2$ the additional score-function terms
appear precisely because the integration over batch-boundary trajectories is replaced by
Monte-Carlo sampling.

\section{The general algorithm}

\begin{algorithm}[H]
\caption{Stochastic Batched DP}
\label{alg:sbdp}
\begin{algorithmic}[1]
\Require Initial $\bc^{(0)}$, total horizon $K_\mathrm{total}$, batch length $K_\mathrm{batch}$, sample size $M$, learning rate $\eta$, iterations $T$
\For{$t = 0, 1, \ldots, T - 1$}
    \State $g \gets 0$
    \For{$m = 1, \ldots, M$ (parallel)}
        \State $\bb \gets p_0$
        \State Initialize trajectory $\tau_m \gets ()$
        \For{$\mathrm{batch} = 1, \ldots, n = K_\mathrm{total} / K_\mathrm{batch}$}
            \State $(V_\mathrm{batch}, \pi_\mathrm{batch}, \mathrm{memo}_\mathrm{batch}) \gets \mathrm{Bellman}(\bb,\, K_\mathrm{batch},\, \bc^{(t)})$
            \For{$k = 1, \ldots, K_\mathrm{batch}$}
                \State $\bs \gets \pi_\mathrm{batch}(\bb)$
                \State $\bx \sim \bb$;  $\by \sim p(\,\cdot \mid \bx, \bs;\, \bc^{(t)})$
                \State $\bb \gets T(\bb, \bs, \by;\, \bc^{(t)})$
                \State Append $(\bs, \by, \bb)$ to $\tau_m$
            \EndFor
        \EndFor
        \State $V_m \gets \Phi_0(\bb;\, \bc^{(t)})$ \Comment{terminal reward at trajectory leaf}
        \State $g_m \gets \mathrm{EnvelopeGrad}(V_m,\, \tau_m,\, \mathrm{memos},\, \bc^{(t)})$
        \State $g \gets g + g_m$
    \EndFor
    \State $g \gets g / M$
    \State $\bc^{(t+1)} \gets \mathrm{Adam}(\bc^{(t)},\, g,\, \eta)$
\EndFor
\State \Return $\bc^{(T)}$
\end{algorithmic}
\end{algorithm}

\textbf{Complexity per outer iteration}: $M \times n_\mathrm{batches} \times
\mathrm{cost}(K_\mathrm{batch} \text{ Bellman})$.  In typical regimes,
$\mathrm{cost}(K_\mathrm{batch} \text{ Bellman}) \approx 10$ seconds, so a single outer
gradient step at $M = 10$ and $n = 4$ takes $\sim 400$ seconds wall-clock without
parallelism, or $\sim 40$ seconds with $M$-way parallelism on a single multicore node.
A complete training run of $T = 500$ iterations finishes in 5--7 hours on standard
hardware.

\section{Convergence and bias}
\label{sec:convergence}

\paragraph{Per-sample exactness.}  Each $V_m$ in Equation~\eqref{eq:estimator} is computed
by exact Bellman within its $K_\mathrm{batch}$ subproblem, with the sharp-max
envelope-theorem gradient.  No smoothing temperature, no policy approximation, no biased
inner optimizer.  The per-sample value and gradient are both exact for the trajectory's
contribution to the estimator.

\paragraph{Estimator bias and variance.}  The estimator
$\widehat{V}^\mathrm{batch}(p_0;\, \bc; \mathcal{M})$ is unbiased for the true batched
value (which is itself an approximation of the full $K_\mathrm{total}$ Bellman value).
Its variance scales as $\sigma^2 / M$ where $\sigma^2$ depends on the spread of
$V_0^{\mathrm{batch}_2}$ across reachable batch-boundary beliefs.  In well-conditioned
problems $\sigma^2$ is small and $M = 10$--$30$ suffices; in problems with bimodal or
multimodal batch-boundary belief distributions $M$ needs to grow accordingly.

\paragraph{Convergence rate.}  Under standard SGD assumptions (bounded gradient variance,
diminishing step sizes), the outer iterates converge to a stationary point of
$V^\mathrm{batch}(p_0;\, \bc)$ at the standard $1/\sqrt{T}$ rate.  Since the per-sample
gradient is exact rather than smoothed, the asymptotic noise floor is set by sampling
variance alone, not by gradient bias.

\paragraph{Approximation gap to full $K_\mathrm{total}$ Bellman.}  The batched policy is a
strict approximation of the full $K_\mathrm{total}$ Bellman policy: each batch's lookahead
is truncated to $K_\mathrm{batch}$, missing the tail value of subsequent batches.  The gap
$V_0(p_0;\, \bc) - V^\mathrm{batch}(p_0;\, \bc) \ge 0$ is regime-dependent.  When the
batch-boundary posterior is unimodal and well-concentrated (e.g., $K_\mathrm{batch}$ is
large enough to "set up" a well-localized belief for subsequent batches), the gap is small.
When the batch-boundary posterior is multimodal, the gap can be large because batched DP
fails to recognize that $\pi^{(1)}$ should anticipate the disambiguation work of $\pi^{(2)}$.

\section{Connections to existing methods}
\label{sec:related}

\paragraph{Stochastic gradient descent.}  The construction is structurally identical to
minibatch SGD on neural networks: full-batch gradient over a finite-but-huge index set
(here, reachable batch-boundary beliefs) is replaced by a small stochastic mini-batch.  The
unbiasedness, variance, and convergence properties carry over directly.

\paragraph{Stochastic Dual Dynamic Programming (SDDP).}  Pereira and Pinto's SDDP for
multi-stage stochastic linear programs uses sampled forward-backward sweeps to build
cutting-plane approximations of the value function.  Stochastic Batched DP is the POMDP
analogue: scenarios are replaced by belief trajectories, cutting-plane approximations are
replaced by exact Bellman within each batch.  The receding-horizon decomposition is a
standard MPC pattern in both communities.

\paragraph{Sequential Bayesian experimental design.}  Simulation-based sequential OED
(Huan and Marzouk; Long, Marzouk, and Wang) uses MC sampling of belief trajectories at
outer optimization time, but typically with rung-3 myopic information-gain inner policies
rather than exact Bellman.  Stochastic Batched DP can be viewed as the upgrade of these
methods from "myopic" to "Bellman within each batch".

\paragraph{Differentiable POMDP planning.}  Existing differentiable-POMDP work (deep RL
with NN policies, Karkus et al.\ QMDP-net, etc.) uses NN-parametrized policies with
end-to-end gradient training.  These methods exhibit a known pathology in joint
hardware-policy optimization: the NN policy can absorb the training signal and leave the
hardware near initialization (the "hardware-light fixed point").  Stochastic Batched DP
sidesteps this by using exact Bellman within each batch, a non-NN inner solver that does
not have a free policy capacity to absorb hardware-task signal.

\paragraph{Sample-based POMDP solvers (POMCP, DESPOT).}  These methods perform Monte-Carlo
tree search over the belief space at policy-inference time.  Stochastic Batched DP uses
sampling at outer-gradient time instead, which is a different point of insertion: the
inner policy is exact Bellman (sample-free), and sampling enters only to estimate the
outer-hardware gradient.

\section{Practical guidance}
\label{sec:practical}

\paragraph{Choosing $K_\mathrm{batch}$.}  Pick the largest $K_\mathrm{batch}$ for which
exact Bellman fits in the available memory budget.  Going larger reduces the
approximation gap to full Bellman; going smaller reduces per-iteration wall-clock at the
cost of suboptimal early-batch lookahead.

\paragraph{Choosing $M$.}  Start at $M = 10$.  Increase if the gradient-step direction is
noisy enough to slow down Adam convergence (visible as oscillation in the loss trajectory).
$M = 30$--$50$ is a typical regime.

\paragraph{Re-sampling cadence.}  Sample fresh trajectories at every outer iteration.
Re-using trajectories across multiple SGD steps introduces gradient bias because the
sampling distribution depends on $\bc^{(t)}$.

\paragraph{Parallelism.}  The $M$ trajectory simulations and Bellman solves are
embarrassingly parallel.  On a $C$-core machine, target $M \le C / k_\mathrm{Bellman}$
where $k_\mathrm{Bellman}$ is the number of cores each individual Bellman solve uses
effectively (typically 8--64).

\paragraph{Diagnostics.}  Track the per-batch values $V_0^{\mathrm{batch}_j}$ and the
spread of $V^{(m)}$ across samples.  A high spread indicates that batch-boundary beliefs
are diverse, which is healthy for sampling but suggests $K_\mathrm{batch}$ may be too
small.  A monotonically decreasing batch-boundary belief variance across batches
indicates progressive disambiguation, which is the regime where batched DP works well.

\section{Discussion}
\label{sec:discussion}

Stochastic Batched DP fills a gap in the relaxation hierarchy between exact Bellman DP
(rung 1, $K_\mathrm{max}$-limited) and the policy-approximation rungs 2--4
(certainty-equivalent, myopic, Gaussian-trace-EKF).  It preserves exactness of the inner
solver within each batch and exactness of per-sample gradients via the envelope theorem,
while breaking the reachable-belief enumeration via Monte-Carlo sampling.  The resulting
method is well-suited to problems where (a) exact Bellman is feasible at modest
$K_\mathrm{batch}$, (b) the deployment horizon is much larger, and (c) the outer hardware
optimization is the central computational expense.

The framework is particularly natural when combined with point-based value iteration
(PBVI) or its rho-POMDP extension as the inner solver.  PBVI's $\alpha$-vector
representation makes the per-sample $K_\mathrm{batch}$ Bellman value $V_0^{\mathrm{batch}_j}(\bb)$
evaluatable at any belief $\bb$ in $O(|\Gamma| \cdot |\mathcal{X}|)$ time without
re-solving from scratch.  Combined with the stochastic-batched outer SGD, this lifts the
horizon ceiling from $K_\mathrm{max} \approx 5$ (exact memoization) to $K_\mathrm{total}$
of practically arbitrary size.

\end{document}
